% GENERATED FILE. Edit canonical lessons, then run npm run generate:books.
% canonical-source-hash: fnv1a64:3fa550e1de53413b
% canonical-lessons: JA-C403-bunka, JA-C403-midokoro, JA-C403-kengaku, JA-C403-tabisaki, JA-C403-otozureru

\chapter{Culture and Places to Visit}
\label{ch:ja-culture-and-places-to-visit}
\hlchaptermodality{403}

\begin{chapteropening}
\textbf{By the end of this chapter:} \emph{I can talk about the culture and highlights of a place.}

Complete the last lesson of chapter 403: I can talk about the culture and highlights of a place.
\end{chapteropening}

\section[\texorpdfstring{\ja{ぶんか} (bunka)}{ぶんか (bunka)}]{\ja{ぶんか} (bunka) — culture}

\label{lesson:JA-C403-bunka}



\begin{quote}
Before the new one: say the Japanese for a local speciality, then the Japanese for to tour.
\end{quote}

\subsection*{You'll want to know: \ja{ぶんか}}
\textbf{\ja{ぶんか}} — \emph{bunka} — \textquotedblleft{}culture\textquotedblright{}.

\begin{quote}
\textbf{\ja{がいこくの} \ja{ぶんかを} \ja{まなびたいです。}} — \emph{gaikoku no bunka o manabitai desu} — \textquotedblleft{}I want to learn about other countries' cultures.\textquotedblright{}
\end{quote}

\begin{cousinweb}
\textbf{\ja{ぶんか}} (\emph{bunka}) is Sino-Japanese (kango): two Chinese characters read with their on'yomi, readings borrowed from Chinese. They are \emph{bun}, \textquotedblleft{}literature\textquotedblright{} (Middle Chinese \emph{mjun}) and \emph{ka}, \textquotedblleft{}change\textquotedblright{} (\emph{xwaeH}) (KANJIDIC2; Middle Chinese in Baxter's transcription).
\end{cousinweb}

\begin{grammarlens}[title={What you've built}]
One more word for getting around.
\end{grammarlens}

\subsection*{Your turn}
\begin{itemize}
\raggedright
  \item \emph{Say it:} \emph{bunka}
  \item \emph{Say it:} \emph{bunka}, once more
  \item \emph{Say it:} \emph{gaikoku no bunka o manabitai desu} — I want to learn about other countries' cultures.
  \item \emph{Recall:} say the Japanese for a local speciality, then the Japanese for to tour, then say \emph{bunka} again
\end{itemize}

\begin{tcolorbox}[breakable,colback=teal!4,colframe=teal!35!black,title={Before you move on}]
What does \ja{ぶんか} mean? (\textquotedblleft{}Culture\textquotedblright{}.) Say it once more.
\end{tcolorbox}
\section[\texorpdfstring{\ja{みどころ} (midokoro)}{みどころ (midokoro)}]{\ja{みどころ} (midokoro) — a highlight; the part worth seeing}

\label{lesson:JA-C403-midokoro}



\begin{quote}
Before the new one: say the Japanese for to tour, then the Japanese for culture.
\end{quote}

\subsection*{You'll want to know: \ja{みどころ}}
\textbf{\ja{みどころ}} — \emph{midokoro} — \textquotedblleft{}a highlight; the part worth seeing\textquotedblright{}.

\begin{quote}
\textbf{\ja{この} \ja{てらの} \ja{みどころは} \ja{にわです。}} — \emph{kono tera no midokoro wa niwa desu} — \textquotedblleft{}The highlight of this temple is its garden.\textquotedblright{}
\end{quote}

\begin{cousinweb}
\textbf{\ja{みどころ}} (\emph{midokoro}) is a native Japanese word (wago), inherited rather than borrowed from Chinese, built from \emph{mi}, \textquotedblleft{}seeing\textquotedblright{} (from \emph{miru}) and \emph{dokoro}, \textquotedblleft{}place\textquotedblright{}; its kanji were attached later as kun'yomi (JMdict).
\end{cousinweb}

\begin{grammarlens}[title={What you've built}]
One more word for getting around.
\end{grammarlens}

\subsection*{Your turn}
\begin{itemize}
\raggedright
  \item \emph{Say it:} \emph{midokoro}
  \item \emph{Say it:} \emph{midokoro}, once more
  \item \emph{Say it:} \emph{kono tera no midokoro wa niwa desu} — The highlight of this temple is its garden.
  \item \emph{Recall:} say the Japanese for to tour, then the Japanese for culture, then say \emph{midokoro} again
\end{itemize}

\begin{tcolorbox}[breakable,colback=teal!4,colframe=teal!35!black,title={Before you move on}]
What does \ja{みどころ} mean? (\textquotedblleft{}A highlight; the part worth seeing\textquotedblright{}.) Say it once more.
\end{tcolorbox}
\section[\texorpdfstring{\ja{けんがく} (kengaku)}{けんがく (kengaku)}]{\ja{けんがく} (kengaku) — a study visit; a tour of a place}

\label{lesson:JA-C403-kengaku}



\begin{quote}
Before the new one: say the Japanese for culture, then the Japanese for a highlight.
\end{quote}

\subsection*{You'll want to know: \ja{けんがく}}
\textbf{\ja{けんがく}} — \emph{kengaku} — \textquotedblleft{}a study visit; a tour of a place\textquotedblright{}.

\begin{quote}
\textbf{\ja{こうじょうを} \ja{けんがくしました。}} — \emph{kōjō o kengaku shimashita} — \textquotedblleft{}We went on a tour of the factory.\textquotedblright{}
\end{quote}

\begin{cousinweb}
\textbf{\ja{けんがく}} (\emph{kengaku}) is Sino-Japanese (kango): two Chinese characters read with their on'yomi, readings borrowed from Chinese. They are \emph{ken}, \textquotedblleft{}see\textquotedblright{} (Middle Chinese \emph{kenH}) and \emph{gaku}, \textquotedblleft{}study\textquotedblright{} (\emph{haewk}) (KANJIDIC2; Middle Chinese in Baxter's transcription).
\end{cousinweb}

\begin{grammarlens}[title={What you've built}]
One more word for getting around.
\end{grammarlens}

\subsection*{Your turn}
\begin{itemize}
\raggedright
  \item \emph{Say it:} \emph{kengaku}
  \item \emph{Say it:} \emph{kengaku}, once more
  \item \emph{Say it:} \emph{kōjō o kengaku shimashita} — We went on a tour of the factory.
  \item \emph{Recall:} say the Japanese for culture, then the Japanese for a highlight, then say \emph{kengaku} again
\end{itemize}

\begin{tcolorbox}[breakable,colback=teal!4,colframe=teal!35!black,title={Before you move on}]
What does \ja{けんがく} mean? (\textquotedblleft{}A study visit; a tour of a place\textquotedblright{}.) Say it once more.
\end{tcolorbox}
\section[\texorpdfstring{\ja{たびさき} (tabisaki)}{たびさき (tabisaki)}]{\ja{たびさき} (tabisaki) — a place visited on a trip}

\label{lesson:JA-C403-tabisaki}



\begin{quote}
Before the new one: say the Japanese for a highlight, then the Japanese for a study visit.
\end{quote}

\subsection*{You'll want to know: \ja{たびさき}}
\textbf{\ja{たびさき}} — \emph{tabisaki} — \textquotedblleft{}a place visited on a trip\textquotedblright{}.

\begin{quote}
\textbf{\ja{たびさきで} \ja{ともだちに} \ja{てがみを} \ja{かきました。}} — \emph{tabisaki de tomodachi ni tegami o kakimashita} — \textquotedblleft{}I wrote to a friend from the place I was visiting.\textquotedblright{}
\end{quote}

\begin{cousinweb}
\textbf{\ja{たびさき}} (\emph{tabisaki}) is a native Japanese word (wago), inherited rather than borrowed from Chinese, built from \emph{tabi}, \textquotedblleft{}journey\textquotedblright{} and \emph{saki}, \textquotedblleft{}ahead\textquotedblright{}; its kanji were attached later as kun'yomi (JMdict).
\end{cousinweb}

\begin{grammarlens}[title={What you've built}]
One more word for getting around.
\end{grammarlens}

\subsection*{Your turn}
\begin{itemize}
\raggedright
  \item \emph{Say it:} \emph{tabisaki}
  \item \emph{Say it:} \emph{tabisaki}, once more
  \item \emph{Say it:} \emph{tabisaki de tomodachi ni tegami o kakimashita} — I wrote to a friend from the place I was visiting.
  \item \emph{Recall:} say the Japanese for a highlight, then the Japanese for a study visit, then say \emph{tabisaki} again
\end{itemize}

\begin{tcolorbox}[breakable,colback=teal!4,colframe=teal!35!black,title={Before you move on}]
What does \ja{たびさき} mean? (\textquotedblleft{}A place visited on a trip\textquotedblright{}.) Say it once more.
\end{tcolorbox}
\section[\texorpdfstring{\ja{おとずれる} (otozureru)}{おとずれる (otozureru)}]{\ja{おとずれる} (otozureru) — to visit a place; to call on}

\label{lesson:JA-C403-otozureru}



\begin{quote}
Before the new one: say the Japanese for a study visit, then the Japanese for a place visited on a trip.
\end{quote}

\subsection*{You'll want to know: \ja{おとずれる}}
\textbf{\ja{おとずれる}} — \emph{otozureru} — \textquotedblleft{}to visit a place; to call on\textquotedblright{}.

\begin{quote}
\textbf{\ja{なつに} \ja{ふるい} \ja{てらを} \ja{おとずれました。}} — \emph{natsu ni furui tera o otozuremashita} — \textquotedblleft{}In summer I visited an old temple.\textquotedblright{}
\end{quote}

Polite \textbf{\ja{おとずれます}}, \emph{otozuremasu}; an ichidan verb. Past (polite): \textbf{\ja{かんこうきゃくが} \ja{しまを} \ja{おとずれました。}}, \emph{kankōkyaku ga shima o otozuremashita}.

\begin{cousinweb}
\textbf{\ja{おとずれる}} (\emph{otozureru}) is a native Japanese word (wago), inherited rather than borrowed from Chinese; its kanji was attached later as a kun'yomi (JMdict).
\end{cousinweb}

\begin{grammarlens}[title={What you've built}]
One more word for getting around.
\end{grammarlens}

\subsection*{Your turn}
\begin{itemize}
\raggedright
  \item \emph{Say it:} \emph{otozureru}
  \item \emph{Say it:} \emph{otozureru}, once more
  \item \emph{Say it:} \emph{natsu ni furui tera o otozuremashita} — In summer I visited an old temple.
  \item \emph{Recall:} say the Japanese for a study visit, then the Japanese for a place visited on a trip, then say \emph{otozureru} again
\end{itemize}

\begin{tcolorbox}[breakable,colback=teal!4,colframe=teal!35!black,title={Before you move on}]
What does \ja{おとずれる} mean? (\textquotedblleft{}To visit a place; to call on\textquotedblright{}.) Say it once more.
\end{tcolorbox}
